\begin{center}
    {\LARGE \textbf{2020无机化学}} \\[2ex]
    {\large 时量180分钟 \quad 满分100分}
\end{center}

\vspace{0.5cm}
\noindent\textbf{注意事项:}
\begin{enumerate}[label=\arabic*., parsep=0pt, itemsep=0pt]
    \item 答题前,考生先将自己的姓名、学号、考场号、座位号填写在答题卡上,将条形码准确粘贴在考生信息条形码粘贴区。
    \item 考试结束后,将本试卷与答题纸一并收回。
    \item 回答各个问题时,将答案写在答题纸上,写在本试卷上无效。
\end{enumerate}

\vspace{0.5cm}
\noindent\textbf{五、鉴别与推断（30 pts.）}

\begin{enumerate}[label=\arabic*., start=1, itemsep=1.5ex]
    \item 选用一种试剂将下列溶液鉴别开，并进一步加以确证。
    \[
    \ce{Na2SO4}, \ce{Na2S2O3}, \ce{Na2SO3}, \ce{Na2S2}, \ce{Na2S}
    \]
    \item 用三种方法鉴别下列各对物质
    \begin{enumerate}[label=(\arabic*), itemsep=1ex]
        \item \ce{SbCl3}和\ce{Bi(NO3)3}
        \item \ce{MnO2}和\ce{PbO2}
    \end{enumerate}
    \item 请回答下列问题：
    \begin{enumerate}[label=(\arabic*), itemsep=1ex]
        \item 在\ce{Pt}制的坩埚底部和壁上粘有高温灼烧过的\ce{Al2O3}，应如何清洗之？
        \item 提纯含有少量铁、钴杂质的金属镍。
        \item 提纯含有少量\ce{Ni^{2+}}杂质的\ce{Co^{2+}}溶液
    \end{enumerate}
    \item 无色晶体\ce{A}溶于稀盐酸得无色溶液，再加入氢氧化钠溶液得白色沉淀\ce{B}，\ce{B}溶于过量氢氧化钠溶液得无色溶液\ce{C}。将\ce{B}溶于盐酸后蒸发、浓缩后又析出\ce{A}。向\ce{A}的稀盐酸溶液加入硫化氢溶液生成橙色沉淀\ce{D}。\ce{D}与过硫化钠可发生氧化还原反应，反应物之间相互作用得到无色溶液\ce{E}。将\ce{A}置于水中生成白色沉淀\ce{F}。写出\ce{A-F}所代表的物质的化学式，并用化学方程式表示\ce{B -> C}和\ce{D -> E}的过程。
    \item 黑蓝色化合物\ce{A}不溶于水，但可溶于稀盐酸生成\ce{B}的溶液，稀释该溶液时，析出白色沉淀\ce{C}，酸化时，沉淀溶解。向\ce{B}的溶液中通硫化氢时，生成灰黑色沉淀\ce{D}，该沉淀可溶于浓盐酸。\ce{D}与硝酸反应时产生浅黄色沉淀\ce{E}，无色气体\ce{F}和化合物\ce{G}的溶液。\ce{F}在空气中迅速变成红棕色。向\ce{G}的溶液中通硫化氢时，生成黄色沉淀\ce{H}。试写出\ce{A}，\ce{B}，\ce{C}，\ce{D}，\ce{G}和\ce{H}所代表的物质的化学式，并写出\ce{D}与硝酸反应的化学方程式。
    \item 将钠盐\ce{A}溶液滴加到\ce{B}溶液（硝酸盐）中，立即有白色沉淀\ce{C}生成，\ce{C}放置一段时间会变为黑色沉淀\ce{D}；若将\ce{B}溶液滴加到\ce{A}溶液中，开始形成的是无色溶液\ce{E}；当加入相当多\ce{B}后才会有\ce{C}沉淀出来。向溶液\ce{B}中加入钾盐\ce{F}溶液，可以生成浅黄色沉淀\ce{G}，\ce{G}可溶于\ce{A}溶液中，\ce{G}在光照下可分解变黑。若向溶液\ce{E}中加入盐酸，有大量沉淀\ce{H}生成，放出刺激性气体\ce{I}。沉淀\ce{H}溶于过量浓硝酸。试写出\ce{A～I}的化学式。
\end{enumerate}







